\section*{Notes on sources}

Questions 1 and 2 come from a lecture board photographed twice \cite{whiteboard1,whiteboard2}; question 3 comes from an assignment known only through a summary \cite{user_statements}, so the claim that the course's ``path'' means no repeated vertex rests on that summary and not on the assignment text, and was not checked against it. The two small marks above $A$ on the board are read as the exponent $k$, as the summary says. The board's drawing has three lines (1--2 with an arrowhead, 2--3, 1--3), while its matrix has two edges, is symmetric, and has a zero diagonal; the paper uses the matrix (\Cref{ex:board}). The board writes the graph as $(G, V)$ and the edge condition as ``$e_{ij}$ in $V$'', and defines connected with ``path'' undefined; the textbook writes $G = (V,E)$ and ``adjacent''.

The textbook contains no walk-counting theorem and no matrix test for connectivity on the pages viewed (printed pp.~11, 12, 34, 61), after a text search of the whole book for the adjacency-matrix phrase; page 7, a notation list, was not viewed, so this absence is not certain. Its authors are given by surname only, and publisher and year were not read. Grinberg and Bai use ``path'' as this paper does \cite{grinberg_graph_theory,bai_bimsa_adjacency}; the textbook's ``path'' is the trail. The scripts are in \texttt{scripts/}: \texttt{m1\_walks.py} and \texttt{m2\_connectivity\_test.py}, run with \texttt{python3 -I} from the paper folder, with their output files beside them.
